\section*{Capítulo \ref{ch:g10:sport-and-health} --- Actividad física y salud}

\begin{solution}{exo:g10:sport-and-health:1}
Un corazón mayor con un volumen sistólico más grande y una frecuencia
de reposo más baja; más capilares y \omterm{def:g10:cells-common-unit:organelle}{mitocondrias} en las \omterm{def:g10:muscles-and-joints:muscle}{fibras
musculares}; más glóbulos rojos; mayores reservas de \omterm{def:g10:exercise-and-energy:glycogen}{glucógeno} (y mejor
regulación del azúcar en sangre y de la presión).
\end{solution}

\begin{solution}{exo:g10:sport-and-health:2}
El volumen sistólico ha crecido, así que los mismos \qty{5}{L/min} se
entregan con menos latidos: $Q = f \times V_s$ con un $V_s$ mayor
necesita un $f$ menor.
\end{solution}

\begin{solution}{exo:g10:sport-and-health:3}
Alrededor de una hora al día de actividad moderada. No: caminar, ir en
bicicleta al instituto, las escaleras y las tareas de casa cuentan
todas --- lo que importa es el movimiento que eleva el \omterm{def:g10:exercise-and-energy:expenditure}{gasto
energético}.
\end{solution}

\begin{solution}{exo:g10:sport-and-health:4}
El uso de sustancias o de métodos que elevan el rendimiento anulando la
regulación del cuerpo en lugar de construir una adaptación. La EPO
anula el control del número de glóbulos rojos; los estimulantes anulan
la fatiga, el aviso que protege al corazón.
\end{solution}

\begin{solution}{exo:g10:sport-and-health:5}
Fatiga persistente, frecuencia cardíaca de reposo en alza y
rendimiento a la baja (también sueño alterado, infecciones frecuentes,
dolor de \omterm{def:g10:muscles-and-joints:muscle}{tendón} o de hueso).
\end{solution}

\begin{solution}{exo:g10:sport-and-health:6}
Al cabo de 12 semanas: $72 - 60 = 12$ latidos con cuatro sesiones y
$72 - 64 = 8$ con dos. La adaptación es proporcional a la demanda ---
más sesiones, más cambio --- y en ambos grupos se frena a medida que se
acerca el nuevo estado.
\end{solution}

\begin{solution}{exo:g10:sport-and-health:7}
$V_s = Q/f$: de $5000/72 \approx \qty{69}{mL}$ a $5000/54 \approx
\qty{93}{mL}$, un factor de $72/54 \approx 1.33$.
\end{solution}

\begin{solution}{exo:g10:sport-and-health:8}
Un $1 - 0.55 = 45\%$ menor. El primer escalón, del quinto menos activo
al siguiente (de 1,00 a 0,80), aporta la mayor ganancia.
\end{solution}

\begin{solution}{exo:g10:sport-and-health:9}
\qty{2.0}{kg}, es decir $2.0/55 \approx 3.6\%$ de su masa. Su volumen
de sangre baja, su frecuencia cardíaca sube para mantener el gasto, su
temperatura corporal trepa y su rendimiento cae; los calambres se
vuelven probables.
\end{solution}

\begin{solution}{exo:g10:sport-and-health:10}
$10 \times 1.1^7 \approx \qty{19.5}{km}$, así que los \qty{20}{km} se
alcanzan en la octava semana: unos dos meses.
\end{solution}

\begin{solution}{exo:g10:sport-and-health:11}
Cada fibra se engrosa añadiendo miofibrillas, y el sistema nervioso
aprende a reclutar más fibras a la vez y mejor sincronizadas; ambas
cosas suben la fuerza. El número de fibras está fijado en la edad
adulta.
\end{solution}

\begin{solution}{exo:g10:sport-and-health:12}
El oxígeno por litro sube $56/45 \approx 1.24$, alrededor de un 24\%;
con el mismo \omterm{def:g10:heart-lungs-effort:output}{gasto cardíaco} y la misma extracción, el
$\dot V\!\mathrm{O_2}$máx sube como mucho esa misma fracción. Pero una
sangre con un 56\% de \omterm{def:g10:cells-common-unit:cell}{células} es mucho más espesa; frenada por un
corazón dormido, puede coagular en un vaso coronario o cerebral.
\end{solution}

\begin{solution}{exo:g10:sport-and-health:13}
En altura, el propio cuerpo sube su EPO en respuesta a una falta real
de oxígeno, en una cantidad controlada, y vuelve a bajarla cuando la
falta termina: la regulación está funcionando. La EPO inyectada impone
una cifra que la regulación nunca elegiría, sin ninguna falta que la
justifique y sin freno.
\end{solution}

\begin{solution}{exo:g10:sport-and-health:14}
El músculo que trabaja capta glucosa sin insulina, y el músculo
entrenado sigue siendo más sensible a ella; caminar baja por tanto el
azúcar en sangre directamente y reduce la insulina necesaria en cada
comida. Un fármaco baja el azúcar pero deja la causa --- un músculo
insensible y sin usar --- donde estaba, y no hace nada por el corazón,
los huesos ni la masa corporal, que caminar también mejora.
\end{solution}

\begin{solution}{exo:g10:sport-and-health:15}
La adaptación se construye en el descanso entre demandas repetidas y
aumentadas poco a poco: un mes de carga intensiva tiene más
probabilidades de producir sobreentrenamiento y lesión que forma
física, ya que el \omterm{def:g10:muscles-and-joints:muscle}{tendón} y el hueso se adaptan en meses, no en semanas.
Y las ganancias se desmontan en unas semanas tras dejarlo: una sola
tanda no ``deja listo'' a nadie para un año. La actividad moderada
regular a lo largo del año hace lo que el mes no puede.
\end{solution}

\begin{solution}{pb:g10:sport-and-health:1}
\textbf{1.} $4800/72 \approx \qty{67}{mL}$ antes; $4800/56 \approx
\qty{86}{mL}$ después.

\textbf{2.} $42 \times 62 \approx \qty{2.6}{L/min}$; $50 \times 62 =
\qty{3.1}{L/min}$: una ganancia del 19\%.

\textbf{3.} Gasto máximo $\dot V\!\mathrm{O_2}/0.150$: de \num{17.4} a
\qty{20.7}{L/min}; volumen sistólico máximo a 200 latidos: de \num{87}
a \qty{103}{mL}.

\textbf{4.} El corazón más grande y más fuerte explica el volumen
sistólico; más capilares y \omterm{def:g10:cells-common-unit:organelle}{mitocondrias} en las fibras explicarían una
ganancia de extracción.

\textbf{5.} Un factor de 4 en 26 semanas es una media de $4^{1/26}
\approx 1.055$, alrededor de un 5,5\% por semana --- dentro de la regla
del 10\% ($1.1^{26}
\approx 12$ habría permitido mucho más).

\textbf{6.} Sobreentrenamiento: frecuencia de reposo en alza, sueño
alterado, infección, rendimiento a la baja y dolor en un tejido de
reparación lenta.

\textbf{7.} Una lesión por sobrecarga del hueso de la espinilla (una
fractura por estrés o una inflamación de su superficie): el hueso es el
tejido que más despacio se adapta, y la carga se duplicó en quince
días.

\textbf{8.} La recuperación entre una sesión y la siguiente es
incompleta; el cuerpo está en un estado de tensión persistente, y la
adaptación que había bajado la frecuencia se está revirtiendo.

\textbf{9.} Una semana de casi reposo y de sueño (la adaptación se
construye durante el descanso); después, dos semanas a la mitad de la
distancia anterior, sin series de velocidad, reconstruyendo una décima
parte por semana; la espinilla vista por un médico y en reposo hasta
que deje de doler (hacer caso al dolor); beber y comer para reparar.

\textbf{10.} Un estimulante: enmascara la fatiga, el aviso que protege
al corazón del agotamiento y del sobrecalentamiento.

\textbf{11.} Sam, $32 \times 80 \approx \qty{2.6}{L/min}$; Álex, $48
\times 64 \approx \qty{3.1}{L/min}$. El de Álex es mayor incluso en
términos absolutos; y la vida diaria --- escaleras, andar, cargar con
uno mismo --- mueve la propia masa del cuerpo, así que la cifra por
kilogramo es la que decide cómo se siente.

\textbf{12.} Sam, $12/32 \approx 38\%$ de su máximo; Álex, $12/48 =
25\%$. Sam, que trabaja por encima de un tercio de su techo, llega sin
aliento.

\textbf{13.} Un músculo sin usar capta poca glucosa de la sangre y
pierde su sensibilidad a la insulina; las mismas comidas necesitan más
insulina cada año y el azúcar en sangre empieza a subir.

\textbf{14.} De 1,00 a 0,70: alrededor de un 30\% menos de riesgo.

\textbf{15.} La \omterm{def:g10:sport-and-health:training}{actividad física} es cualquier movimiento que eleve el
\omterm{def:g10:exercise-and-energy:expenditure}{gasto energético}: ir andando al trabajo, las escaleras en vez del
ascensor, las tareas de casa --- media hora al día de eso no necesita
ni deporte ni equipamiento.

\textbf{16.} Al 1\% anual, un margen del 10\% son diez años de
descenso.

\textbf{17.} El máximo de Sam fue menor y lleva veinticinco años
perdiendo hueso; su muñeca ha cruzado el umbral en el que una caída la
rompe, mientras que el hueso de Álex, construido bajo carga en la
adolescencia y mantenido desde entonces, tiene diez años de margen.

\textbf{18.} Una adaptación es una estructura que se mantiene por el
uso, no una adquisición permanente: se desmonta cuando la demanda cesa
y se reconstruye cuando vuelve.

\textbf{19.} Enfermedades del corazón y de las arterias, diabetes de
tipo 2 y osteoporosis (también algunos cánceres y la depresión).
Diabetes: el músculo inactivo pierde su sensibilidad a la insulina, el
azúcar en sangre sube y el páncreas se sobrecarga.

\textbf{20.} Frecuencia cardíaca de reposo 54 frente a 78;
$\dot V\!\mathrm{O_2}$máx \num{48} frente a \qty{32}{mL/min} por
kilogramo; unas tres horas de carrera a la semana, mantenidas durante
catorce años.
\end{solution}
